% Part: incompleteness
% Chapter: representability-in-q
% Section: beta-function

\documentclass[../../../include/open-logic-section]{subfiles}

\begin{document}

\olfileid{inc}{req}{bet}
\olsection{د بېټا تابعې لمه}


د دې د ښودلو دپاره چې که جمع، ضرب، او~$\Char{=}$ راسره وي نو
بنسټيز بازګښت ترسره کولے شو، بايد داسې تابعې جوړې کړو چې لړۍ سمبال
کړي۔ (که توان هم راسره وے، کار به مو اسان و۔) کله چې بنسټيز بازګښت
راسره و، د «$n$م اول عدد» په شان څيزونه مو تعريفولے شول او نسبتاً
ساده کوډونه مو غوره کولے شول۔ خو دلته بنسټيز بازګښت راسره نشته---په
حقيقت کښې غواړو وښيو چې بنسټيز بازګښت د اقل موندلو په کارولو ترسره
کولے شو---نو بايد څه زياته ځيرکتيا وکاروو۔

\begin{lem}
\ollabel{lem:beta}
داسې تابعه~$\beta(d,i)$ شته چې د هرې لړۍ $a_0$، \dots،~$a_n$
دپاره داسې عدد~$d$ شته چې د هر~$i \le n$ دپاره
$\beta(d,i) = a_i$ وي۔ سربېره پر دې،~$\beta$ له بنسټيزو تابعو څخه
يوازې د ترکيب او منظم اقل موندلو په کارولو تعريفېدے شي۔
\end{lem}

$d$ داسې وګڼئ چې لړۍ~$\tuple{a_0, \dots, a_n}$ کوډوي، او
$\beta(d,i)$ ئې~$i$م غړے بېرته ورکوي۔ (پام وکړئ چې دا «کوډول» هغه
د اول عددونو د توانونو کوډول \emph{نۀ} کاروي چې لا له وړاندې ورسره
اشنا يو!) لمه ډېره لږه ادعا کوي؛ دا نۀ وايي چې لړۍ يو ځاے کولے شو،
غړي ورزياتولے شو، يا ان دا چې د ترکيب او منظم اقل موندلو له لارې
تعريفېدونکو تابعو په کارولو له~$a_0$، \dots،~$a_n$ څخه~$d$
\emph{محاسبه} کولے شو۔ يوازې دومره وايي چې د «کوډ پرانيستلو» داسې
تابعه شته چې هره لړۍ «کوډ شوې» وي۔

د~$\beta$ نښه د ګوډل ده۔ بيا وايو، د لمې د ثبوت ګرانه برخه دا ده
چې د ظاهراً محدودو وسيلو، يعنې يوازې د ترکيب او اقل موندلو، په
کارولو مناسبه~$\beta$ تعريف کړو---خو اجازه لرو چې جمع، ضرب،
او~$\Char{=}$ وکاروو۔ د دې لمې د ثبوت بېلابېلې لارې شته، خو يوه له
تر ټولو پاکو لارو څخه لا هم د ګوډل اصلي لاره ده، چې د عددونو د تيورۍ
هغه حقيقت ئې کاراوه چې د سونزي قضيه بلل کېږي (په دوديز ډول، «د چين
د پاتې شونو قضيه»)۔

\begin{defn}
دوه طبيعي عددونه~$a$ او~$b$ هله \emph{خپلمنځي اول} دي چې ستر
ګډ مقسوم‌عليه ئې~$1$ وي؛ په بله وينا، له دې پرته بل ګډ مقسوم‌عليه
نۀ لري۔
\end{defn}

\begin{defn}
طبيعي عددونه~$a$ او~$b$ \emph{د~$c$ په مودولو سره مطابق} دي،
$a \equiv b \mod c$، هله چې~$c \mid (a-b)$؛ يعنې کله چې پر~$c$
ووېشل شي،~$a$ او~$b$ يو شان پاتې شونې لري۔
\end{defn}

دا د سونزي قضيه ده:
\begin{thm}
فرض کړئ $x_0$، \dots،~$x_n$ (جوړه په جوړه) خپلمنځي اول دي۔
$y_0$، \dots،~$y_n$ هر عددونه وي۔ بيا داسې عدد~$z$ شته چې
\begin{align*}
z & \equiv y_0 \mod x_0 \\
z & \equiv y_1 \mod x_1 \\
& \vdots  \\
z & \equiv y_n \mod x_n.
\end{align*}
\end{thm}

د سونزي قضيه به داسې کاروو: که $x_0$، \dots،~$x_n$ په ترتيب سره
له~$y_0$، \dots،~$y_n$ څخه لوی وي،~$z$ د لړۍ
$\tuple{y_0, \dots,y_n}$ کوډ ګڼلے شو۔ د~$y_i$ د بېرته موندلو
دپاره يوازې~$z$ پر~$x_i$ وېشو او پاتې شونې ئې اخلو۔ د دې کوډ کارولو
دپاره بايد د~$x_0$، \dots،~$x_n$ مناسب ارزښتونه ومومو۔

په دې برخه کښې به څو کتنې راسره مرسته وکړي۔ د ورکړل شوو
$y_0$، \dots،~$y_n$ دپاره دا ونيسئ:
\begin{align*}
j &= \max(n, y_0 + 1, \dots, y_n + 1), \\
m &= \lcm(1,\dots,j),
\end{align*}
او دا ونيسئ:
\begin{align*}
x_0 & = 1 + m \\
x_1 & = 1 + 2 \cdot m \\
x_2 & = 1 + 3 \cdot m \\
& \vdots  \\
x_n & = 1 + (n+1) \cdot m
\end{align*}
بيا دوه خبرې سمې دي:
\begin{enumerate}
\item\ollabel{rel-prime} $x_0,\dots,x_n$ خپلمنځي اول دي۔
\item\ollabel{less} د هر~$i$ دپاره، $y_i < x_i$۔
\end{enumerate}
د دې ليدلو دپاره چې~\olref{rel-prime} سمه ده، پام وکړئ چې که~$p$
اول عدد وي او $p \mid x_i$ او~$p \mid x_k$ وي، نو
$p \mid 1 + (i+1) m$ او~$p \mid 1 + (k+1) m$ دي۔ خو بيا~$p$ د
دواړو توپير وېشي:
\[
(1 + (i+1)m) - (1+ (k+1)m) = (i-k) m.
\]
څرنګه چې~$p$، $1 + (i+1)m$ وېشي، په عين وخت کښې~$m$ نشي وېشلے
(که ئې وېشلے، لومړۍ وېشنې به~$1$ پاتې شونې پرېښې وه)۔ نو ځکه چې
$p$، $(i-k)m$ وېشي،~$p$ بايد~$i-k$ ووېشي۔ خو~$\left|i-k\right|$ تر
ډېره~$n$ دے، او موږ~$j \geq n$ غوره کړے دے، نو له دې څخه بيا
$p \mid m$ لازمېږي، چې تناقض دے۔ نو هېڅ اول عدد~$x_i$ او~$x_k$
دواړه نۀ وېشي۔ فقره~\olref{less} اسانه ده:
$y_i < j \leq m < x_i$ لرو۔

اوس د~$\beta$ تابعې لمه ثابته کړو۔ په ياد ولرئ چې~$0$، تالي،
جمع، ضرب،~$\Char{=}$، پروجکشنونه، او هره هغه تابعه کارولے شو چې
له دوى څخه د ترکيب او پر منظمو تابعو د اقل موندلو په کارولو تعريف
شوې وي۔ که د کومې اړيکې ځانګړونکې تابعه همداسې تعريفېدونکې وي، هغه
اړيکه هم کارولے شو۔ د پخوا په شان ښودلے شو چې دا اړيکې د بولي
ترکيبونو او محدود کمیت ايښودنې په وړاندې تړلې دي؛ د بېلګې په توګه:
\begin{align*}
\fn{not}(x) & \defis \Char{=}(x,0)\\
\bmin{x \leq z}{R(x,y)} & \defis \umin{x}{(R(x,y) \lor x = z)}\\
\bexists{x \leq z}{R(x,y)} & \defiff R(\bmin{x \leq z}{R(x,y)}, y)
\end{align*}
بيا ښودلے شو چې دا ټول هم له بنسټيز بازګښت پرته تعريفېدونکي دي:
\begin{enumerate}
\item جوړوونکې تابعه، $J(x,y) = \frac{1}{2}[(x+y)(x+y+1)] + x$؛
% maybe explain more what is going on here, a bit confusing.
\item د پروجکشن تابعې
\begin{align*}
K(z) & = \bmin{x \leq z}{\bexists{y \leq z}{z = J(x,y)}},\\
L(z) & = \bmin{y \leq z}{\bexists{x \leq z}{z = J(x,y)}};
\end{align*}
\item د کوچني والي اړيکه~$x < y$؛
\item د وېش اړيکه~$x \mid y$؛
% \item $x \tsub y$
% \item $\fn{Prime}(x)$
% \item Assuming $p$ is prime, the relation ``$x$ is a power of $p$'':
% \[
% \bforall{y \leq x}{(y \mid x \lif y = 1 \lor y = x)}.
% \]
\item تابعه~$\fn{rem}(x,y)$، چې د~$y$ پر~$x$ د وېش پاتې شونې
  بېرته ورکوي۔
\end{enumerate}
اوس تعريف کړئ:
\begin{align*}
\beta^*(d_0,d_1,i) & = \fn{rem}(1+(i+1) d_1,d_0) \text{ او}\\
\beta(d,i) & = \beta^*(K(d),L(d),i).
\end{align*}
دا هماغه تابعه ده چې غواړو ئې۔ پورته غوندې ورکړل شوو
$a_0,\dots,a_n$ دپاره دا ونيسئ:
\[
j = \max(n,a_0+1,\dots,a_n+1),
\]
او~$d_1 = \lcm(1,\dots,j)$ ونيسئ۔ له پورته
\olref{rel-prime} څخه پوهېږو چې $1+d_1$، $1+2 d_1$، \dots،
$1+(n+1) d_1$ خپلمنځي اول دي، او له~\olref{less} څخه چې دا ټول
له~$a_0,\dots,a_n$ څخه لوی دي۔ د سونزي د قضيې له مخې داسې ارزښت
$d_0$ شته چې د هر~$i$ دپاره
\[
d_0 \equiv a_i \mod (1+(i+1)d_1)
\]
وي، او نو (ځکه~$d_1$ له~$a_i$ څخه لوی دے)،
\[
a_i = \fn{rem}(1+(i+1)d_1,d_0).
\]
$d = J(d_0,d_1)$ ونيسئ۔ بيا د هر~$i \le n$ دپاره لرو:
\begin{align*}
\beta(d,i) & =  \beta^*(d_0,d_1,i) \\
& =  \fn{rem}(1+(i+1) d_1,d_0) \\
& =  a_i
\end{align*}
او دا هماغه څه دي چې پکار مو وو۔ په دې سره د~$\beta$ تابعې د لمې
ثبوت بشپړېږي۔

\begin{prob}
  وښيئ چې اړيکې~$x < y$، $x \mid y$، او تابعه
  $\fn{rem}(x,y)$ له بنسټيز بازګښت پرته تعريفېدے شي۔ $0$، تالي،
  جمع، ضرب،~$\Char{=}$، پروجکشنونه، او محدوده کمينه‌موندنه او کمیت
  ايښودنه کارولے شئ۔
\end{prob}

\end{document}
